\documentclass[aspectratio=169,10pt]{beamer}
\makeatletter\def\input@path{{../estilo/}{estilo/}}\makeatother
\usepackage{devmobile}
\modulo{08}{Rede e Consumo de APIs}{HTTP, REST e JSON, Retrofit, OkHttp, tratamento de erros, offline-first, imagens e segurança}

\begin{document}
\capa

\begin{frame}{Objetivos do módulo}
  \begin{conceito}[Ao final deste módulo você será capaz de]
    \begin{itemize}
      \item Descrever o ciclo requisição/resposta do HTTP, seus métodos e códigos de status.
      \item Ler a documentação de uma API REST e modelar seus dados em Kotlin.
      \item Consumir APIs com \textbf{Retrofit} + \textbf{kotlinx.serialization} + \textbf{OkHttp}.
      \item Tratar falhas de rede e representá-las no estado da tela.
      \item Implementar a estratégia \emph{offline-first} com Room como fonte da verdade.
      \item Carregar imagens e proteger chaves e dados em trânsito.
    \end{itemize}
  \end{conceito}
\end{frame}

\sumario

% =============================================================================
\section{Fundamentos de HTTP}

\begin{frame}[fragile]{Requisição e resposta}
  \begin{center}
  \begin{tikzpicture}[node distance=30mm]
    \node (cel) {\celular[0.5]{\tiny app}};
    \node[caixaescura,minimum width=2.6cm,minimum height=16mm,right=40mm of cel] (srv) {Servidor\\ \tiny api.exemplo.com};
    \draw[seta] ([yshift=5mm]cel.east) -- node[above,font=\scriptsize,align=center]{\textbf{Requisição}\\ \cod{GET /bancadas/3} + cabeçalhos} ([yshift=5mm]srv.west);
    \draw[setaambar] ([yshift=-5mm]srv.west) -- node[below,font=\scriptsize,align=center]{\textbf{Resposta}\\ \cod{200 OK} + JSON} ([yshift=-5mm]cel.east);
  \end{tikzpicture}
  \end{center}
  \begin{columns}[T]
    \begin{column}{0.49\textwidth}
      \begin{codigo}[title=Requisição]{bash}
GET /v1/bancadas/3 HTTP/1.1
Host: api.labmonitor.ifma.edu.br
Accept: application/json
Authorization: Bearer eyJhbGciOi...
      \end{codigo}
    \end{column}
    \begin{column}{0.49\textwidth}
      \begin{codigo}[title=Resposta]{JSON}
HTTP/1.1 200 OK
Content-Type: application/json

{ "id": 3, "nome": "Bancada 3",
  "temp_c": 24.8, "livre": true }
      \end{codigo}
    \end{column}
  \end{columns}
\end{frame}

\begin{frame}{Métodos e códigos de status}
  \begin{columns}[T]
    \begin{column}{0.46\textwidth}
      \footnotesize
      \begin{tabularx}{\linewidth}{@{}l L@{}}
        \toprule
        \cabecalho{Método} & \cabecalho{Semântica} \\
        \midrule
        \cod{GET} & ler um recurso (sem corpo) \\
        \cod{POST} & criar um recurso ou disparar ação \\
        \cod{PUT} & substituir um recurso inteiro \\
        \cod{PATCH} & alterar parte de um recurso \\
        \cod{DELETE} & remover um recurso \\
        \bottomrule
      \end{tabularx}
      \begin{conceito}[Idempotência]
        \cod{GET}, \cod{PUT} e \cod{DELETE} podem ser repetidos sem efeito extra, por isso é seguro tentar de novo após uma falha. \cod{POST}, em geral, não.
      \end{conceito}
    \end{column}
    \begin{column}{0.52\textwidth}
      \footnotesize
      \begin{tabularx}{\linewidth}{@{}l L@{}}
        \toprule
        \cabecalho{Faixa} & \cabecalho{Significado e exemplos} \\
        \midrule
        \textcolor{verdeok}{\textbf{2xx}} & sucesso: \cod{200 OK}, \cod{201 Created}, \cod{204 No Content} \\
        \textbf{3xx} & redirecionamento: \cod{301}, \cod{304 Not Modified} \\
        \textcolor{ambar}{\textbf{4xx}} & erro do cliente: \cod{400} inválido, \cod{401} sem login, \cod{403} sem permissão, \cod{404} não existe, \cod{409} conflito, \cod{429} limite de requisições \\
        \textcolor{vermelhoerro}{\textbf{5xx}} & erro do servidor: \cod{500}, \cod{502}, \cod{503} indisponível \\
        \bottomrule
      \end{tabularx}
      \begin{dica}
        4xx: o app precisa corrigir algo (login, dados). 5xx e falhas de conexão: vale tentar de novo mais tarde.
      \end{dica}
    \end{column}
  \end{columns}
\end{frame}

% =============================================================================
\section{APIs REST e JSON}

\begin{frame}{Projetando a API do LabMonitor}
  \footnotesize
  \begin{tabularx}{\textwidth}{@{}l l L@{}}
    \toprule
    \cabecalho{Método} & \cabecalho{Endpoint} & \cabecalho{Ação} \\
    \midrule
    \cod{GET} & \cod{/v1/bancadas} & lista todas as bancadas \\
    \cod{GET} & \cod{/v1/bancadas/\{id\}} & detalha uma bancada \\
    \cod{GET} & \cod{/v1/bancadas/\{id\}/leituras?desde=...} & leituras a partir de uma data (query string) \\
    \cod{POST} & \cod{/v1/reservas} & cria uma reserva (corpo JSON) \\
    \cod{DELETE} & \cod{/v1/reservas/\{id\}} & cancela uma reserva \\
    \bottomrule
  \end{tabularx}
  \vspace{4pt}
  \begin{columns}[T]
    \begin{column}{0.48\textwidth}
      \begin{conceito}[REST em uma frase]
        URLs identificam \textbf{recursos} (substantivos), métodos HTTP dizem \textbf{o que fazer} (verbos), e o estado trafega em JSON.
      \end{conceito}
    \end{column}
    \begin{column}{0.48\textwidth}
      \begin{dica}[API pública para praticar]
        A \textbf{Open-Meteo} não exige chave:\\
        {\scriptsize\url{https://api.open-meteo.com/v1/forecast?latitude=-2.53&longitude=-44.30&current=temperature_2m}}\\
        devolve a temperatura atual em São Luís.
      \end{dica}
    \end{column}
  \end{columns}
\end{frame}

% =============================================================================
\section{Configuração no Android}

\begin{frame}[fragile]{Permissão e HTTPS}
  \begin{columns}[T]
    \begin{column}{0.5\textwidth}
\begin{codigo}[title=AndroidManifest.xml]{XML}
<uses-permission
  android:name="android.permission.INTERNET" />

<application
  android:networkSecurityConfig=
    "@xml/network_security_config" ... >
\end{codigo}
\begin{codigo}[title=res/xml/network\_security\_config.xml]{XML}
<network-security-config>
  <!-- HTTP liberado SÓ para o ESP32 na rede local -->
  <domain-config cleartextTrafficPermitted="true">
    <domain includeSubdomains="false">192.168.4.1</domain>
  </domain-config>
</network-security-config>
\end{codigo}
    \end{column}
    \begin{column}{0.48\textwidth}
      \begin{conceito}[HTTPS por padrão]
        Desde o Android 9, o tráfego em texto puro (\cod{http://}) é bloqueado. Use \cod{https://} sempre que possível.
      \end{conceito}
      \begin{atencao}[Dispositivos IoT]
        Placas como o ESP32 muitas vezes servem HTTP simples na rede local. Libere \textbf{apenas aquele endereço}, nunca \cod{cleartextTrafficPermitted="true"} global.
      \end{atencao}
      \begin{dica}
        No emulador, \cod{10.0.2.2} aponta para o \cod{localhost} do seu computador.
      \end{dica}
    \end{column}
  \end{columns}
\end{frame}

% =============================================================================
\section{Retrofit e kotlinx.serialization}

\begin{frame}[fragile]{Modelando o JSON}
  \begin{columns}[T]
    \begin{column}{0.5\textwidth}
\begin{codigo}[title=Resposta da Open-Meteo (resumo)]{JSON}
{
  "latitude": -2.53,
  "longitude": -44.3,
  "current": {
    "time": "2026-11-01T15:00",
    "temperature_2m": 30.4,
    "relative_humidity_2m": 62
  }
}
\end{codigo}
      \begin{dica}
        Ferramentas como o \emph{JSON to Kotlin Class} (plugin da IDE) geram os DTOs a partir de um exemplo.
      \end{dica}
    \end{column}
    \begin{column}{0.48\textwidth}
\begin{codigo}[title=ClimaDto.kt]{Kotlin}
@Serializable
data class ClimaDto(
  val latitude: Double,
  val longitude: Double,
  val current: AtualDto
)

@Serializable
data class AtualDto(
  val time: String,
  @SerialName("temperature_2m")
  val temperatura: Double,
  @SerialName("relative_humidity_2m")
  val umidade: Int? = null     // opcional
)
\end{codigo}
      \begin{conceito}
        \cod{@SerialName} liga o nome do JSON ao nome Kotlin. Valores padrão tornam campos opcionais.
      \end{conceito}
    \end{column}
  \end{columns}
\end{frame}

\begin{frame}[fragile]{A interface da API}
  \begin{columns}[T]
    \begin{column}{0.54\textwidth}
\begin{codigo}[title=LabApi.kt]{Kotlin}
interface LabApi {
  @GET("v1/bancadas")
  suspend fun listarBancadas(): List<BancadaDto>

  @GET("v1/bancadas/{id}")
  suspend fun bancada(@Path("id") id: Int): BancadaDto

  @GET("v1/bancadas/{id}/leituras")
  suspend fun leituras(
    @Path("id") id: Int,
    @Query("desde") desde: String
  ): List<LeituraDto>

  @POST("v1/reservas")
  suspend fun reservar(@Body r: NovaReservaDto): ReservaDto

  @DELETE("v1/reservas/{id}")
  suspend fun cancelar(@Path("id") id: Long): Response<Unit>
}
\end{codigo}
    \end{column}
    \begin{column}{0.44\textwidth}
      \scriptsize
      \begin{tabularx}{\linewidth}{@{}l L@{}}
        \toprule
        \cabecalho{Anotação} & \cabecalho{Função} \\
        \midrule
        \cod{@GET}, \cod{@POST}\ldots & método e caminho \\
        \cod{@Path} & substitui \cod{\{id\}} na URL \\
        \cod{@Query} & \cod{?chave=valor} \\
        \cod{@Body} & objeto serializado no corpo \\
        \cod{@Header} & cabeçalho específico \\
        \bottomrule
      \end{tabularx}
      \begin{conceito}[Funções \cod{suspend}]
        O Retrofit executa a chamada fora da thread principal e devolve o objeto já convertido. Status fora de 2xx lança \cod{HttpException}.
      \end{conceito}
      \begin{dica}
        Use \cod{Response<T>} quando precisar ler o código de status ou os cabeçalhos.
      \end{dica}
    \end{column}
  \end{columns}
\end{frame}

\begin{frame}[fragile]{Montando o cliente com OkHttp e Hilt}
  \begin{columns}[T]
    \begin{column}{0.58\textwidth}
\begin{codigo}[title=di/NetworkModule.kt]{Kotlin}
@Module @InstallIn(SingletonComponent::class)
object NetworkModule {
  @Provides @Singleton
  fun okHttp(sessao: SessaoRepository) = OkHttpClient.Builder()
    .connectTimeout(10, TimeUnit.SECONDS)
    .addInterceptor { chain ->         // token em toda requisição
      chain.proceed(chain.request().newBuilder()
        .header("Authorization", "Bearer ${sessao.token()}")
        .build())
    }
    .addInterceptor(HttpLoggingInterceptor().apply {
      level = if (BuildConfig.DEBUG) Level.BODY else Level.NONE
    }).build()
  @Provides @Singleton
  fun api(client: OkHttpClient): LabApi = Retrofit.Builder()
    .baseUrl("https://api.labmonitor.ifma.edu.br/")
    .client(client)
    .addConverterFactory(Json { ignoreUnknownKeys = true }
      .asConverterFactory("application/json".toMediaType()))
    .build().create(LabApi::class.java)
}
\end{codigo}
    \end{column}
    \begin{column}{0.4\textwidth}
      \begin{conceito}[Camadas]
        \textbf{Retrofit} transforma a interface em chamadas HTTP; \textbf{OkHttp} faz a conexão (cache, timeouts, interceptadores); \textbf{kotlinx.serialization} converte JSON.
      \end{conceito}
      \begin{dica}[\cod{ignoreUnknownKeys}]
        Se a API ganhar campos novos, o app antigo continua funcionando.
      \end{dica}
      \begin{atencao}
        Log de corpo (\cod{BODY}) só em debug: ele mostra tokens e dados pessoais.
      \end{atencao}
    \end{column}
  \end{columns}
\end{frame}

% =============================================================================
\section{Erros e estados}

\begin{frame}[fragile]{O que pode dar errado}
  \begin{columns}[T]
    \begin{column}{0.5\textwidth}
      \scriptsize
      \begin{tabularx}{\linewidth}{@{}l L@{}}
        \toprule
        \cabecalho{Exceção} & \cabecalho{Causa provável} \\
        \midrule
        \cod{UnknownHostException} & sem internet / DNS \\
        \cod{SocketTimeoutException} & servidor lento \\
        \cod{HttpException} (4xx/5xx) & resposta de erro \\
        \cod{SerializationException} & JSON diferente do esperado \\
        \cod{SSLHandshakeException} & certificado inválido \\
        \bottomrule
      \end{tabularx}
      \begin{atencao}
        Todas as de rede herdam de \cod{IOException}. Trate-as \textbf{no repositório} e entregue à UI um erro com significado.
      \end{atencao}
    \end{column}
    \begin{column}{0.48\textwidth}
\begin{codigo}[title=Erros com significado]{Kotlin}
sealed interface ErroRede {
  data object SemConexao : ErroRede
  data object NaoAutorizado : ErroRede
  data class Servidor(val codigo: Int) : ErroRede
  data object Desconhecido : ErroRede
}

suspend fun <T> chamarApi(
  bloco: suspend () -> T
): Result<T> = try {
  Result.success(bloco())
} catch (e: CancellationException) { throw e
} catch (e: IOException) {
  Result.failure(FalhaRede(ErroRede.SemConexao))
} catch (e: HttpException) {
  Result.failure(FalhaRede(when (e.code()) {
    401, 403 -> ErroRede.NaoAutorizado
    else -> ErroRede.Servidor(e.code()) }))
}
\end{codigo}
    \end{column}
  \end{columns}
\end{frame}

\begin{frame}[fragile]{Do erro à tela}
  \begin{columns}[T]
    \begin{column}{0.54\textwidth}
\begin{codigo}[title=ClimaViewModel.kt]{Kotlin}
fun carregar() = viewModelScope.launch {
  _estado.update { it.copy(carregando = true, erro = null) }
  chamarApi { api.clima(-2.53, -44.30) }
    .onSuccess { dto ->
      _estado.update {
        it.copy(carregando = false,
          temperatura = dto.current.temperatura)
      }
    }
    .onFailure { e ->
      _estado.update {
        it.copy(carregando = false,
          erro = (e as? FalhaRede)?.tipo?.mensagem()
            ?: "Erro inesperado")
      }
    }
}
\end{codigo}
    \end{column}
    \begin{column}{0.44\textwidth}
      \centering
      \begin{tikzpicture}[node distance=3mm]
        \node[caixa,minimum width=4.2cm] (c) {\textbf{Carregando}: indicador de progresso};
        \node[caixaescura,minimum width=4.2cm,below=of c] (s) {\textbf{Sucesso}: dados na tela};
        \node[caixaambar,minimum width=4.2cm,below=of s] (v) {\textbf{Vazio}: mensagem + ação};
        \node[caixa,draw=vermelhoerro,fill=vermelhoerro!8,minimum width=4.2cm,below=of v] (e) {\textbf{Erro}: explicação + ``Tentar de novo''};
      \end{tikzpicture}
      \begin{dica}[Repetir com espera crescente]
        Para falhas temporárias, tente de novo após 1\,s, 2\,s, 4\,s\ldots (\emph{exponential backoff}). O operador \cod{retryWhen} do Flow e o WorkManager fazem isso.
      \end{dica}
    \end{column}
  \end{columns}
\end{frame}

% =============================================================================
\section{Offline-first}

\begin{frame}{Estratégia offline-first}
  \begin{columns}[c]
    \begin{column}{0.52\textwidth}
      \centering
      \begin{tikzpicture}[node distance=7mm]
        \node[caixaescura,minimum width=2.6cm] (ui) {UI / ViewModel};
        \node[caixapetroleo,minimum width=2.6cm,below=of ui] (repo) {Repositório};
        \node[caixaambar,minimum width=2.2cm,below left=7mm and -6mm of repo] (db) {Room\\ \tiny fonte da verdade};
        \node[caixa,minimum width=2.2cm,below right=7mm and -6mm of repo] (api) {API remota};
        \node[caixa,minimum width=2.2cm,below=8mm of api,fill=cinzaclaro] (wm) {WorkManager\\ \tiny sincronização};
        \draw[seta] (db) -- node[rotulo,left]{\cod{Flow}} (repo);
        \draw[seta] (repo) -- node[rotulo,right]{\cod{StateFlow}} (ui);
        \draw[setaambar] (api.west) -- node[rotulo,below]{grava} (db.east);
        \draw[setatrac] (wm) -- node[rotulo,right]{agenda} (api);
        \draw[setatrac] (db.south) |- node[rotulo,below,pos=0.75]{fila de envios} (wm.west);
      \end{tikzpicture}
    \end{column}
    \begin{column}{0.46\textwidth}
      \begin{conceito}[Leitura]
        A UI lê \textbf{sempre} do banco local. A rede apenas atualiza o banco, e a tela reage sozinha.
      \end{conceito}
      \begin{conceito}[Escrita]
        Grave localmente na hora (marcado como ``pendente'') e envie ao servidor quando houver conexão, com WorkManager (Módulo 10).
      \end{conceito}
      \begin{dica}[Resultado]
        O app abre instantaneamente, funciona no laboratório sem Wi-Fi e economiza dados móveis.
      \end{dica}
    \end{column}
  \end{columns}
\end{frame}

\begin{frame}[fragile]{Implementando no repositório}
  \begin{columns}[T]
    \begin{column}{0.56\textwidth}
\begin{codigo}[title=BancadaRepositoryImpl.kt]{Kotlin}
override fun bancadas(): Flow<List<Bancada>> =
  dao.observarTodas()
    .map { lista -> lista.map { it.paraModelo() } }
    .onStart { atualizarSeAntigo() }   // dispara refresh

private suspend fun atualizarSeAntigo() {
  if (agora() - ajustes.ultimaSync() < 5.minutes) return
  chamarApi { api.listarBancadas() }.onSuccess { dtos ->
    dao.substituirTodas(dtos.map { it.paraEntidade() })
    ajustes.registrarSync(agora())
  } // em falha: a UI segue com os dados em cache
}

override suspend fun reservar(r: NovaReserva) {
  dao.inserirReserva(r.paraEntidade(pendente = true))
  agendador.agendarEnvioPendentes()   // WorkManager
}
\end{codigo}
    \end{column}
    \begin{column}{0.42\textwidth}
      \begin{atencao}[Conflitos]
        Se duas pessoas reservarem a mesma bancada offline, o servidor decide. Trate \cod{409 Conflict} avisando o usuário e removendo a reserva local.
      \end{atencao}
      \begin{dica}[Mostre a idade dos dados]
        ``Atualizado há 3 min'' e um ícone de nuvem riscado quando offline deixam claro o que o usuário está vendo.
      \end{dica}
      \begin{dica}
        Monitore a conectividade com o \cod{ConnectivityManager} (\cod{registerDefault}\-\cod{NetworkCallback}) convertido em \cod{Flow} via \cod{callbackFlow}.
      \end{dica}
    \end{column}
  \end{columns}
\end{frame}

% =============================================================================
\section{Imagens, alternativas e segurança}

\begin{frame}[fragile]{Imagens com Coil e alternativas ao Retrofit}
  \begin{columns}[T]
    \begin{column}{0.5\textwidth}
\begin{codigo}[title=Coil em Compose]{Kotlin}
AsyncImage(
  model = ImageRequest.Builder(LocalContext.current)
    .data(equipamento.fotoUrl)
    .crossfade(true)
    .build(),
  placeholder = painterResource(R.drawable.carregando),
  error = painterResource(R.drawable.sem_foto),
  contentDescription = equipamento.descricao,
  contentScale = ContentScale.Crop,
  modifier = Modifier.size(72.dp)
    .clip(RoundedCornerShape(8.dp))
)
\end{codigo}
      \begin{conceito}[Coil]
        Baixa, redimensiona e guarda imagens em cache (memória e disco) automaticamente.
      \end{conceito}
    \end{column}
    \begin{column}{0.48\textwidth}
      \scriptsize
      \begin{tabularx}{\linewidth}{@{}l L@{}}
        \toprule
        \cabecalho{Biblioteca} & \cabecalho{Quando considerar} \\
        \midrule
        \textbf{Retrofit} & padrão de mercado em apps Android \\
        \textbf{Ktor Client} & projetos Kotlin Multiplatform (Módulo 13) \\
        \textbf{OkHttp puro} & chamadas avulsas, download de arquivos \\
        \textbf{WebSocket} (OkHttp/Ktor) & dados em tempo real bidirecionais \\
        \textbf{MQTT} & telemetria IoT (Módulo 10) \\
        \bottomrule
      \end{tabularx}
      \begin{dica}
        A arquitetura em camadas permite trocar Retrofit por Ktor mexendo só na fonte remota.
      \end{dica}
    \end{column}
  \end{columns}
\end{frame}

\begin{frame}{Segurança na comunicação}
  \begin{columns}[T]
    \begin{column}{0.48\textwidth}
      \begin{erro}[Nunca coloque segredos no app]
        Tudo o que vai no APK pode ser extraído: chaves de API, senhas de banco, tokens de serviços pagos. \textbf{Ofuscação não protege segredos.}
      \end{erro}
      \begin{conceito}[Solução: backend intermediário]
        O app chama \textbf{o seu servidor} (ou uma Cloud Function), que guarda a chave e chama o serviço externo. O app se autentica com o token do usuário.
      \end{conceito}
    \end{column}
    \begin{column}{0.5\textwidth}
      \centering
      \begin{tikzpicture}[node distance=7mm]
        \node (a) {\celular[0.35]{}};
        \node[caixaescura,right=of a] (b) {Seu backend\\ \tiny guarda a chave};
        \node[caixa,right=of b] (c) {Serviço\\ externo};
        \draw[seta] (a) -- node[rotulo,above]{token do usuário} (b);
        \draw[seta] (b) -- node[rotulo,above]{chave secreta} (c);
      \end{tikzpicture}
      \begin{dica}[Checklist]
        \begin{itemize}
          \item HTTPS em tudo.
          \item Tokens de curta duração, com renovação.
          \item Valide tudo também no servidor.
          \item Limite de requisições no backend.
          \item Logs de rede desligados em release.
        \end{itemize}
      \end{dica}
    \end{column}
  \end{columns}
\end{frame}

% =============================================================================
\section{Fechamento}

\begin{frame}{Resumo}
  \begin{resumo}
    \begin{itemize}
      \item HTTP = requisição (método, URL, cabeçalhos, corpo) + resposta (status, cabeçalhos, corpo).
      \item DTOs \cod{@Serializable} espelham o JSON; Retrofit transforma a interface em chamadas \cod{suspend}.
      \item OkHttp cuida de timeouts, interceptadores (token, log) e cache.
      \item Converta exceções em erros com significado no repositório e mostre estados claros na UI.
      \item Offline-first: a UI lê do Room, a rede atualiza o banco e as escritas entram numa fila.
      \item Nada de segredos no app: use um backend intermediário e HTTPS.
    \end{itemize}
  \end{resumo}
\end{frame}

\begin{frame}{Exercícios}
  \begin{exercicio}[Exercício 1 --- Open-Meteo]
    Crie uma tela que mostra temperatura e umidade atuais de São Luís usando a Open-Meteo, com os estados carregando, sucesso e erro, e um botão ``Atualizar''.
  \end{exercicio}
  \begin{exercicio}[Exercício 2 --- Falhas]
    Teste o app em modo avião, com \emph{throttling} de rede no emulador e com uma URL inexistente. Cada caso mostra uma mensagem adequada?
  \end{exercicio}
  \begin{exercicio}[Exercício 3 --- ESP32 local]
    Programe um ESP32 para servir \cod{GET /leitura} com JSON na rede local. Consuma-o no app com Retrofit, liberando HTTP só para o IP da placa.
  \end{exercicio}
  \begin{projeto}[LabMonitor --- Etapa 8]
    Adicione a fonte remota (API real ou um servidor \emph{mock}) ao repositório de bancadas com estratégia offline-first, mostrando ``atualizado há X min'' no painel.
  \end{projeto}
\end{frame}

\begin{frame}{Para aprofundar}
  \small
  \begin{itemize}
    \item Conectar à rede: \url{https://developer.android.com/develop/connectivity/network-ops/connecting}
    \item Network security config: \url{https://developer.android.com/privacy-and-security/security-config}
    \item Retrofit: \url{https://square.github.io/retrofit/}
    \item kotlinx.serialization: \url{https://github.com/Kotlin/kotlinx.serialization}
    \item Apps offline-first: \url{https://developer.android.com/topic/architecture/data-layer/offline-first}
    \item Coil: \url{https://coil-kt.github.io/coil/} \quad Open-Meteo: \url{https://open-meteo.com}
  \end{itemize}
\end{frame}

\encerramento{Módulo 09 --- Firebase: autenticação, Firestore, regras de segurança, armazenamento e notificações push.}

\end{document}
